\section{A direct-count approximation interface}
\label{sec:approximate-proofs}

This section changes the representation of the half-gate statistics, not
its statistical mechanism.  An earlier twelve-coordinate construction
already established the real-error half-gate argument and the finite-bank
precision, noise, and source-edit phenomenon \citep{BCJ001}.  Here the supplied
coordinates directly approximate $J,C,D$, rather than reconstructing $J$
from two nonnegative successor coordinates.  Sequential clipping makes
this changed interface legal.  In particular, the result is not a claim
that exact recovery of a discontinuous gate bit is necessary, nor a new
precision principle.

\subsection{Statement and clipping}

For every row $u=(a,b)$, retain $d(u)=1-a$, and define the clean counts
\begin{align}
 J_u&=\sum_{i=0}^{M-1}\ind\{Y_i=u\},&
 C_u&=\sum_{i=0}^{M-1}\ind\{Y_i=u,\ X_{i+2}=d(u)\},
 \label{eq:a-full-counts}\\
 D_u&=\sum_{i=0}^{\lfloor M/2\rfloor-1}\ind\{Y_i=u\}.
 \label{eq:a-half-count}
\end{align}
Thus $D$ denotes a first-half departure count, always normalized by $M$
when used as a frequency.  Primed counts are evaluated on the corrupted
word $Z$.  At the actual protected row $u=Y_M$, omit the subscript and
write $q=\theta_u(d(u))$.  The protected pair makes the selected row the
same for the clean and corrupted words.

For real $x$ and $a\le b$, write
$\operatorname{clip}(x;a,b)=\min\{b,\max\{a,x\}\}$.
Suppose an observation-measurable computation supplies, for every row,
real values $\widetilde J_u,\widetilde C_u,\widetilde D_u$.  Its only
accuracy premise is the simultaneous pathwise bound
\begin{equation}\label{eq:a-stat-error}
 \max_u\max\{
 |\widetilde J_u-J'_u|,
 |\widetilde C_u-C'_u|,
 |\widetilde D_u-D'_u|\}\le e,
 \qquad e\ge0.
\end{equation}
This must hold on every allowed word and for every outcome of the
computation's independent random coins.  The computation has no access
to the clean word or hidden source parameters.  Its errors may be
dependent across coordinates and may depend on the observed word.
The number $e$ is a proved error bound, not an input required by the
decoder.  Clip sequentially, using the already clipped first coordinate:
\begin{equation}\label{eq:a-sequential-clip}
 \overline J_u=\operatorname{clip}(\widetilde J_u;0,M),\qquad
 \overline C_u=\operatorname{clip}(\widetilde C_u;0,\overline J_u),\qquad
 \overline D_u=\operatorname{clip}(\widetilde D_u;0,\overline J_u).
\end{equation}
The approximate half-gate returns designated probability
\begin{equation}\label{eq:a-gate}
 \widehat q=
 \begin{cases}
 \overline C/\overline J,
       &\overline J>0\ \text{and}\ 512\overline D\ge\overline J,\\
 0,&\text{otherwise}.
 \end{cases}
\end{equation}
Equality accepts.  A final observation-measurable probability
$\widehat q_{\rm out}\in[0,1]$ with
$|\widehat q_{\rm out}-\widehat q|\le\xi$, $\xi\ge0$, is also allowed.
The complementary symbol receives probability $1-\widehat q_{\rm out}$.

\begin{theorem}[Approximate direct half-gate statistics]
\label{thm:approximate-gate}
For every integer $M\ge64$, every integer $0\le k\le M$, every homogeneous
binary order-two source with an all-state certificate
$P^2(u,v)\ge\beta\Pi(v)$, $\beta>0$, every initial law, and every legal
source-independent edit kernel in \eqref{eq:attack}, the decoder
\eqref{eq:a-sequential-clip}--\eqref{eq:a-gate}, with the pathwise premise
\eqref{eq:a-stat-error}, has unsquared total-variation risk
\begin{equation}\label{eq:a-risk-general}
 R\le\min\left\{1,
 \frac{65536}{\sqrt M}+\frac{65536}{M\sqrt\beta}
 +1072176\frac{6k+e}{M}+\xi\right\}.
\end{equation}
The same decoder works simultaneously for all these parameters without
knowing $k,\beta$, or $e$.  At the supplied certificate $\beta=1/M$,
\begin{equation}\label{eq:a-risk-critical}
 R\le\min\left\{1,
 \frac{131072}{\sqrt M}+6433056\frac{k}{M}
 +1072176\frac eM+\xi\right\}.
\end{equation}
The edit kernel does not observe the predictor's random seed, and the
result permits arbitrary independent randomization satisfying the
pathwise accuracy premise.
\end{theorem}

\begin{lemma}[Clipping with a perturbed upper endpoint]
\label{lem:a-clipping}
Let $0\le c\le j\le M$ and
$|\widetilde j-j|,|\widetilde c-c|\le e$.
Set $\overline j=\operatorname{clip}(\widetilde j;0,M)$ and
$\overline c=\operatorname{clip}(\widetilde c;0,\overline j)$.
Then $|\overline j-j|\le e$ and $|\overline c-c|\le e$.
The same statement holds with $c$ replaced by any $d\in[0,j]$.
\end{lemma}
\begin{proof}
Projection onto $[0,M]$ does not increase distance to $j$.
For the second coordinate,
$\overline c\le\max\{\widetilde c,0\}\le c+e$.
For its lower bound, if $\widetilde c<0$, then $c\le e$, so
$\overline c=0\ge c-e$.  If
$0\le\widetilde c\le\overline j$, clipping retains
$\widetilde c\ge c-e$.  If $\widetilde c>\overline j$, then
$\overline c=\overline j\ge j-e\ge c-e$.
These cases exhaust the possibilities, including $\overline j=0$;
no clipped coordinate is required to be integer.
\end{proof}

\subsection{Complete pathwise perturbation argument}

One symbol edit affects at most three departure triples.  Changing a
triple has histogram $\ell_1$ cost at most two, so each full-count
coordinate changes by at most $6k$.  Restricting to the first half does
not increase this bound.  Adjacent edits and triples crossing the split
are covered by the union bound; their affected triples are not assumed
disjoint.  Lemma~\ref{lem:a-clipping} therefore gives, on every edit and
computation outcome,
\begin{equation}\label{eq:a-clean-errors}
 |\overline J-J|,
 |\overline C-C|,
 |\overline D-D|\le\Delta:=6k+e.
\end{equation}
The perturbation radius $\Delta$ may be an arbitrarily small noninteger.
Define the events
\[
 G=\{\overline J>0,\ 512\overline D\ge\overline J\},\qquad
 H=\{D\ge J/1024\},\qquad
 B=\{J>0,\ D<J/256\}.
\]
Both $C/J$, where defined, and the accepted clipped ratio belong to
$[0,1]$.

On $G\cap H\cap\{J>0\}$, we claim
\begin{equation}\label{eq:a-ratio}
 \left|\frac{\overline C}{\overline J}-q\right|
 \le\left|\frac CJ-q\right|+\frac{8\Delta}{J+1}.
\end{equation}
If $\Delta=0$, the ratios coincide.  If $\Delta>0$ and $J\ge2\Delta$,
then $\overline J\ge J/2$, and subtracting ratios gives
\[
 \left|\frac{\overline C}{\overline J}-\frac CJ\right|
 \le\frac{|\overline C-C|}{\overline J}
       +\frac{C|J-\overline J|}{J\overline J}
 \le\frac{4\Delta}{J}
 \le\frac{8\Delta}{J+1},
\]
where the last step uses the integer inequality $J\ge1$.
If $0<J<2\Delta$, then $\Delta>J/2$ and
$8\Delta/(J+1)\ge1$, so the unit bound on the ratio difference suffices.
This proves \eqref{eq:a-ratio} for all real $\Delta\ge0$.

On $G\cap H^c$, acceptance and \eqref{eq:a-clean-errors} imply
\[
 D+\Delta\ge\overline D
 \ge\frac{\overline J}{512}\ge\frac{J-\Delta}{512},
 \qquad D<\frac J{1024}.
\]
Combining them yields $J<1026\Delta$.
On $G^c\cap B^c\cap\{J>0\}$, if $\overline J=0$ then
$J\le\Delta$.  Otherwise the strict rejection inequality and
$D\ge J/256$ give
\[
 \frac J{256}-\Delta
 \le\overline D<\frac{\overline J}{512}
 \le\frac{J+\Delta}{512},
\]
and hence $J<513\Delta$.  For $\Delta=0$, neither false-status case
with $J>0$ is possible.  Rejection on $B$ has loss exactly $q$.

The occupation lemma has an integer threshold, so introduce
\begin{equation}\label{eq:a-ceilings}
 w_1=\max\{1,\lceil1026\Delta\rceil\},\qquad
 w_2=\max\{1,\lceil513\Delta\rceil\}.
\end{equation}
They cover the strict false-status inequalities and $J=0$.
In the separate $\overline J=0$, $J>0$ case, $J\le\Delta$ forces
$\Delta>0$, and then $J<513\Delta$ as required.
For every $\Delta\ge0$,
$w_1+w_2\le1539\Delta+2$, including $\Delta=0$.
The preceding cases, bounding any remaining loss by one, prove the
pointwise reduction
\begin{equation}\label{eq:a-pathwise}
 \begin{split}
 |\widehat q-q|\le{}&
 \ind\{J>0\}\left|\frac CJ-q\right|
 +8\Delta\frac{\ind_H}{J+1}+q\ind_B\\
 &+\ind\{J<w_1\}+\ind\{J<w_2\}.
 \end{split}
\end{equation}
The ratio $C/J$ is evaluated only when $J>0$.  At $J=0$ the tail
indicators cover the loss without defining that ratio.  All right-hand
terms are nonnegative.  In particular, no consistency assumption that
the approximate coordinates are counts of a word is used, and no exact
agreement with the corrupted word's gate bit is asserted.

\subsection{Averaging and normalized precision}

\begin{proof}[Proof of Theorem~\ref{thm:approximate-gate}]
The four stochastic inputs have complete proofs in
Section~\ref{sec:positive-proofs}.  With $D=J_{1,Y_M}$, they are precisely
\begin{align*}
 \Prb(J<w)&\le\min\{1,16w/(M+1)\}
          &&\text{for integer }w\ge1
          &&\text{(Lemma~\ref{lem:occupation})},\\
 \E\!\left[\ind\{J>0\}\left|C/J-q\right|\right]
      &\le512/\sqrt{M+1}
          &&&&\text{(Lemma~\ref{lem:empirical})},\\
 \E[q\ind_B]&\le65536/(M\sqrt\beta)
          &&&&\text{(Lemma~\ref{lem:imbalance})},\\
 \E[\ind_H/(J+1)]&\le130944/M
          &&&&\text{(\eqref{eq:p-half-inverse} in
                    Proposition~\ref{prop:finite-upper})}.
\end{align*}
In particular these inputs hold for arbitrary starts, including starts
outside stationary support.  They require neither endpoint-conditioned
independent departures nor independence of the perturbation errors.
Apply them to \eqref{eq:a-pathwise}, then use \eqref{eq:a-ceilings}:
\[
 \begin{split}
 \E|\widehat q-q|\le{}&
 \frac{512}{\sqrt{M+1}}+\frac{65536}{M\sqrt\beta}
 +\frac{8\cdot130944\Delta}{M}
 +\frac{16(1539\Delta+2)}{M+1}\\
 \le{}&\frac{65536}{\sqrt M}+\frac{65536}{M\sqrt\beta}
       +\frac{1072176\Delta}{M}.
 \end{split}
\]
Here
\[
 8\cdot130944+16(1026+513)=1072176,\qquad
 \frac{512}{\sqrt{M+1}}+\frac{32}{M+1}
 \le\frac{65536}{\sqrt M}.
\]
The final output error contributes at most $\xi$ by the triangle
inequality.  A legal binary probability has loss at most one.
Substitute $\Delta=6k+e$ to prove \eqref{eq:a-risk-general};
then $\beta=1/M$ proves \eqref{eq:a-risk-critical}.
The proof is pointwise before averaging over all permissible
randomizations, which proves the stated simultaneous quantifiers.
\end{proof}

For a nonnegative integer $h$, let $Q=2^h$ and separately round each
of the twelve normalized primed statistics
$J'_u/M,C'_u/M,D'_u/M$ to its nearest point of
$\{0,1/Q,\ldots,1\}$, with any fixed tie rule.  Multiplication by $M$
gives approximations satisfying $e\le M/(2Q)$.  Thus
\eqref{eq:a-risk-critical} gives the explicit rounded-bank upper bound
\begin{equation}\label{eq:a-rounded-risk}
 R\le\min\left\{1,
 \frac{131072}{\sqrt M}+6433056\frac kM+\frac{536088}{Q}+\xi
 \right\}.
\end{equation}
Consequently $h\ge\lceil\tfrac12\log_2 M\rceil$ and
$\xi\le M^{-1/2}$ suffice for a constant multiple of $g_M(k)$,
uniformly in the unknown budget.  If a separate output stage has error
at most $2^{-b}$, its contribution is simply $\xi=2^{-b}$; the output
precision $b$ is distinct from the bank precision $h$ and the raw-window
width $r$.

For known $M$, a straightforward integer-numerator representation of
this bank uses at most $12(h+1)$ bits, plus the two exact protected
endpoint bits.  The extra numerator bit permits the grid value one.
This is a representation statement, not a streaming algorithm with
that workspace: rounding after each update can accumulate error.
Nor does it prove that all twelve coordinates are necessary.
Raw input storage, exact prefix selection, addressing, accumulation,
pair/triple feature construction, division, output precision, weights,
depth, attention work, and training are separate costs.  In particular,
a classical counter has the same statistical interface; there is no
neural statistical advantage in \eqref{eq:a-rounded-risk}.
